%%%%%%%%%%  scan idx 188  |  folio 175  %%%%%%%%%%
sont pleinement fidèles et ont pour image essentielle $\mathrm D^b(S)_{\coh}$ (I 2.15). En outre, si $S$ est de «~dimension cohomologique finie~» au sens de (3.7)\footnote{Le \S~3 est indépendant du reste de l'exposé.} (par exemple si l'espace sous-jacent à $S$ est noethérien et de dimension de Krull finie), alors les foncteurs
\[
\mathrm K^{-}(L(S))/\mathrm K^{-,\varnothing}(L(S))\longrightarrow
\mathrm K^{-}(\Loclib(S))/\mathrm K^{-,\varnothing}(\Loclib(S))\longrightarrow
\mathrm D^{-}(S)_{\coh}
\]
sont des équivalences de catégories.
\item Le foncteur canonique
\[
\mathrm K^b(\Loclib(S))/\mathrm K^{b,\varnothing}(\Loclib(S))\longrightarrow\mathrm D(S)
\]
est pleinement fidèle et a pour image essentielle $\mathrm D(S)_{\parf}$ (I 4.9). Pour que $E\in\Ob\mathrm D(S)$ soit d'amplitude parfaite contenue dans $[m,n]$ (I 4.8), il faut et il suffit que $E$ soit isomorphe, dans $\mathrm D(S)$, à un complexe $F$ tel que $F^i=0$ pour $i\notin[m,n]$, $F^i\in\Ob L(S)$ pour $i>m$, $F^m$ étant localement libre de type fini.
\end{enumerate}

\emph{Preuve.} Compte tenu de (3.5) et (3.7)\footnotemark[\value{footnote}], c'est un corollaire immédiat de (2.2.8)~: le seul point à vérifier, qui résulte d'ailleurs trivialement de (3.5) (resp. (3.7)), est que pour un objet $E$ de $\mathrm D^b(\Qcoh(S))$ (resp. $\mathrm D(\Qcoh(S))$ dans le cas où $S$ est de dimension cohomologique finie au sens de (3.7)), les notions de pseudo-cohérence et de perfection de (2.2.1) coïncident avec les notions usuelles (I 2.15 et 4.9) lorsqu'on regarde $E$ comme objet de $\mathrm D(S)$.

\textbf{Remarque 2.2.10.} a) Si $S$ est un schéma affine d'anneau $A$, la famille réduite au faisceau $\mathcal O_S$ est ample. Pour celle-ci, on a $L(S)=\Lib(S)$. Il est clair par ailleurs que les catégories $\mathrm K^{-,\varnothing}(\Lib(S))$, et $\mathrm K^{-,\varnothing}(\Loclib(S))$ sont réduites aux objets nuls. Par suite, on conclut de (2.2.9) que, si l'on note $\Mod(A)$ la catégorie des $A$-modules et $\mathrm D(A)=\mathrm D(\Mod(A))$, le foncteur $\mathbb R\Gamma_S:\mathrm D^+(S)\longrightarrow\mathrm D^+(A)$ induit un foncteur pleinement fidèle $\mathrm D^b(S)_{\coh}\longrightarrow\mathrm D(A)$ (resp. $\mathrm D(S)_{\parf}\longrightarrow\mathrm D(A)$) dont l'image essentielle est 
